\chapter*{Preface}


\section*{June 2025}

We organised a series of performance analysis workshops in spring 2025, funded through the Digital Research Infrastructure (DRI) programme under the umbrella of our HAI-End project. 
These workshops differed from previous installments in that we did not focus on particular tools. 
Instead, we conducted brainstorming and discussion exercises centered around a fundamental question: 
How do you begin performance assessment when you know little to nothing about your code?


While some tools such as Intel's Application Performance Snapshots (APS) or Linaro's MAP provide high-level overviews, we approached this challenge not from a tools perspective but from an observational point of view. 
We asked: 
What are the essential properties we must understand before diving deeper into code analysis?


This inquiry led to a high-level, top-down approach, documented through mindmaps, flowcharts, and case studies. 
The raising approach is fundamentally different from the existing POP methodology, which collects fine-grained metrics and successively aggregates them into high-level metrics and flaws. 
It also differs from approaches offered by and supported through tools like Scalasca, which require a certain level of code maturity and system knowledge from the user.


Both direction of travel are important and very valuable, and we think they might be more sophisticated and insightful compared to what we came up with.
However, our technique is very simple and does require next to no specialised software and tools knowledge.
We developed a straightforward, accessible starting point to initiating performance analysis.
This first version of the document summarises these foundational ideas.


Some of the sections are not yet written, and a lot more are  yet to come. 
However, we decided to share them early, to be able to gather feedback and polish the ideas.
This is only the first version of a living document.


Take home: 
- Target group has to be clearer: in-depth or application-domain focused
- Tool information density plus training material is not appropriate for latter group
- Future workshops have to cover both domains separately

- Application people want to do performance engineering but don't know how to do (and that they have to do) assessment first
- They pick the first tool and then stick to it and hence often miss out on big picture

- Success model: pair up tool expert with domain expert
- Tool training is there for assessor, but not necessarily for the domain expert who wants to do performance engineering
- This contradicts my previous assumption that tools are there predominantly for performance engineer

- Maybe the assessments have to come first
- Success stories are key (like the ones POP has delivered)

- A lot of tools are too difficult to use; counterpart is MAQAO or APS or Perf Snapshot

- Pitch assessment gurus directly for tool workshops
- Create for others a tutorial "how to assess your code's performance" 


\begin{flushright}
Thomas Flynn,
Tobias Weinzierl
\\
Durham, June 6, 2025
\end{flushright}


